\runninghead{Opening and Map}
\section{Waiting When God Seems Silent}

A new Church year begins without fanfare. The \work{Gloria in excelsis} falls silent under the Roman Missal's rubric for Advent Sundays (GIRM~53), and the first voice heard from the ambo is an accusation mixed with pleading. In the closing chapters of the book of Isaiah, a bruised people asks God why he has let them wander from his ways and hardened their hearts so that they no longer fear him, begs him to tear open the heavens and come down, confesses that every righteous deed they can show has grown stained, and ends by placing itself in his hands like wet clay before a potter. Psalm~80(79) takes up a kindred cry: the Shepherd of Israel is asked to wake his hidden strength, turn his people back, and let his face shine upon a vineyard that has suffered ruin. Then the tone shifts sharply. St.~Paul addresses the Christians of Corinth with overflowing gratitude, declaring that they lack no spiritual gift as they wait for the unveiling of Jesus Christ, and assuring them that a faithful God will keep them blameless to the end. At last, on the Mount of Olives, Christ gives a parting parable of a householder who travels abroad, entrusts his servants with authority over every task, stations a doorkeeper, and leaves everyone under a single charge: \emph{Watch}, for no one knows in which watch of the night the master of the house will return.

On the First Sunday of Advent in Year~B (celebrated in the United States on Sunday, November~29, 2026), those four passages set a sharp contrast before the assembly. How can the same liturgy begin by admitting that no one rises up to take hold of God (Isaiah~64:7 in the Vulgate; 64:6 in Hebrew numbering), move to Paul's confidence that believers are already enriched in Christ and sustained by God's fidelity (1~Corinthians~1:4--9), and conclude with Mark's urgent warning against being caught asleep when the master arrives suddenly (Mark~13:33--37)? The postconciliar Mass places prayers shared across all three Sunday cycles of Advent beside a Year~B course of readings from the \work{Ordo lectionum Missae}. Ancient and later Catholic interpreters worked closely through each of those passages, often reaching genuine disagreements over divine hardening, human frailty, Christ's withheld knowledge of the final hour, and the relation between grace and moral watchfulness.

Examining each appointed passage in its own setting before weighing three interpretations grounded in the Fathers and Doctors shows how much turns on which exegetical difficulty stands in the foreground. Readers who wish to orient each scriptural passage in time and place may turn to the chronological table in the \hyperref[sec:date-location]{Scriptural Date and Location} appendix, which separates traditional attributions, historical settings, and scholarly composition proposals. Because modern copyrighted English liturgical translations and protected modern Latin orations cannot be reproduced here, all English Scripture quotations use the public-domain Douay--Rheims (Challoner) text clearly marked as study English, while modern Missal prayers are identified by their Latin opening words and liturgical citations. The table below maps all twelve appointed parts of the Mass in liturgical order.

\begin{maptable}
Entrance Antiphon & \emph{Ad te levavi} (Ps.~25[24]:1--3) & Lifts the soul to God in trust that those who wait on him will not be put to shame. \\
Collect & \emph{Da, quaesumus, omnipotens Deus} & Opening prayer of the First Sunday of Advent (modern text withheld; its ancient Gelasian antecedent asks God for the will to meet the coming Christ in righteous deeds). \\
First Reading & Isa.~63:16b--17, 19b; 64:2b--7 (Vg.~63:16b--17; 64:1, 3b--8) & Confesses wandering, sin, and stained righteousness; begs God to tear open the heavens; appeals to him as Father and Potter over clay. \\
Responsorial Psalm & Ps.~80(79):2ac, 3b, 15--16, 18--19; R.~v.~4 & Pleads for the Shepherd enthroned upon the Cherubim to stir up his might, shine his face, visit his vineyard, and confirm the Son of man. \\
Second Reading & 1~Cor.~1:3--9 & Thanks God for the grace enriching the Corinthians in speech and knowledge as they wait for Christ's unveiling; promises that faithful God calls them into his Son's fellowship. \\
Gospel Acclamation & \emph{Ostende nobis, Domine} (Ps.~85[84]:8) & Framed by \emph{Alleluia}; asks the Lord to show his mercy and grant his salvation. \\
Gospel & Mark~13:33--37 & Charges the disciples and all believers to stay alert and pray like servants and a doorkeeper awaiting a master who may return at any watch of the night. \\
Prayer over the Offerings & \emph{Suscipe, quaesumus, Domine} & Prayer over the bread and wine (modern text withheld; its ancient Veronense antecedent offers gifts drawn from God's own benefits toward eternal redemption). \\
Seasonal Preface & Advent Preface & Seasonal Eucharistic thanksgiving of Advent (two Prefaces of Advent stand in the \work{Ordo Missae}; modern text withheld). \\
Communion Antiphon & \emph{Dominus dabit benignitatem} (Ps.~85[84]:13) & Declares that the Lord will give goodness (or sweetness) and that our earth will yield its fruit. \\
Prayer after Communion & \emph{Prosint nobis} & Closing prayer of the Communion Rite (modern text withheld). \\
Concluding Blessing & Ordinary blessing, or optional solemn blessing / prayer over the people & Closes the Mass before the dismissal under GIRM~90 (modern solemn-blessing texts withheld). \\
\end{maptable}
